% =====================================================================
%  CAPÍTULO 8 — CONCLUSÃO
% =====================================================================

\chapter{CONCLUSÃO}
\label{cap:conclusao}

Este trabalho apresentou o \argus{}, um ecossistema de software multiplataforma
e de código aberto para simulação, exploração e análise de percursos reais no
ciclismo virtual. O ponto de partida não foi a construção de mais um catálogo de
rotas, mas a constatação de que o percurso e o dado fisiológico costumam ocupar
aplicativos diferentes. No \textit{desktop}, o artefato foi projetado para
manter o GPX, a simulação, a telemetria e o histórico no mesmo fluxo, com
processamento local e fórmulas explícitas; no Android, oferece uma sessão ao
vivo com o mesmo \textit{frontend}, sem fechar o ciclo de FIT e análise.

\section{Síntese do percurso}
\label{sec:sintese-final}

O trabalho percorreu as etapas da \textit{Design Science Research}: identificou
o problema da fragmentação, definiu requisitos, projetou e desenvolveu o
artefato, demonstrou o fluxo em cenários de reconhecimento, exploração e treino,
e avaliou o resultado em um estudo piloto com o próprio autor. A versão
registrada no INPI foi tratada como linha de base; a versão apresentada incorpora
o modo Android, a migração do mapa, a análise fisiológica, as integrações e o
assistente local.

A tese do trabalho se confirma no plano da engenharia: é possível reunir esses
elementos em um único ambiente, com uma arquitetura que separa domínio,
interface e persistência, e com uma interface reutilizada entre \textit{desktop}
e Android. A confirmação é limitada pelo escopo: o estudo piloto não demonstra
eficácia de treinamento, usabilidade por terceiros ou superioridade sobre
plataformas comerciais. Demonstra que o fluxo funciona nas condições testadas e
que as decisões de projeto são viáveis.

\section{Contribuições}
\label{sec:contribuicoes}

As contribuições do trabalho podem ser organizadas em cinco grupos.

\begin{enumerate}
    \item \textbf{Artefato aberto e integrável}: o \argus{} é distribuído sob
        GPLv3 e, no \textit{desktop}, reúne importação de GPX, simulação de
        percurso, controle de rolo, telemetria, análise fisiológica e
        histórico local em um único sistema.

    \item \textbf{Análise fisiológica transparente}: no \textit{desktop}, as
        métricas NP, IF, TSS, TRIMP, PMC, HRR e desacoplamento são calculadas
        com fórmulas e parâmetros documentados, e o texto explicita o que é
        medido, o que é estimado e o que é convenção de treino.

    \item \textbf{Arquitetura multiplataforma delimitada}: o sistema foi
        exercitado no \textit{desktop} Windows/Linux (fluxo completo) e no
        telefone Android (sessão ao vivo), com uma preparação arquitetural para
        iOS/iPadOS e um layout responsivo para tablet. O trabalho não transforma
        preparação em validação.

    \item \textbf{Processo documentado}: o histórico Git, os ciclos de
        desenvolvimento, os requisitos e as limitações permitem que outro
        pesquisador entenda e reproduza o percurso de construção.

    \item \textbf{Registro e continuidade}: o registro no INPI formaliza a
        versão inicial, e o texto documenta a evolução até a versão atual,
        preservando a memória do projeto.
\end{enumerate}

\section{Limitações}
\label{sec:limitacoes-finais}

As limitações são parte do resultado, não um apêndice incômodo. As principais
são:

\begin{itemize}
    \item \textbf{avaliação com um único participante}: o piloto foi realizado
        pelo autor, sem outros voluntários, sem grupo de controle e sem
        instrumento padronizado de usabilidade;
    \item \textbf{cobertura de plataformas}: iOS/iPadOS, tablet Android e
        \textit{desktop} macOS não foram validados em dispositivo;
    \item \textbf{escopo funcional do Android}: o modo móvel cobre a sessão ao
        vivo, mas não exporta FIT, não mantém histórico e não inclui testes de
        aptidão, rolo virtual ou integrações;
    \item \textbf{cobertura de hardware}: o piloto usou um conjunto restrito de
        rolo, sensor de frequência cardíaca e demais sensores; outros fabricantes
        e protocolos podem revelar problemas de compatibilidade;
    \item \textbf{métricas fisiológicas}: VT1/VT2 são marcados manualmente, o
        teste de cinco minutos é usado como proxy de MAP e não como medida de
        VO\textsubscript{2}max, o TRIMP usa um coeficiente populacional e as
        zonas de FC são uma aproximação por percentual da máxima;
    \item \textbf{reconhecimento não automatizado}: o fluxo de reconhecimento é
        manual; não há detecção automática de segmentos, curvas ou
        recomendações de \textit{pacing};
    \item \textbf{dependência de tiles}: o mapa 3D usa tiles remotos; no
        \textit{desktop}, o processamento e o histórico são locais, mas a
        visualização completa do mapa não é totalmente offline;
    \item \textbf{ausência de multiplayer e de catálogo}: o sistema não oferece
        corrida em tempo real com outros atletas nem biblioteca própria de
        rotas.
\end{itemize}

\section{Trabalhos futuros}
\label{sec:trabalhos-futuros}

A agenda de continuidade tem dois eixos: aprofundar a avaliação e ampliar o
artefato.

No eixo de avaliação, os próximos passos são:

\begin{itemize}
    \item realizar um estudo de usabilidade com ciclistas amadores, com
        instrumento padronizado e registro de tarefas;
    \item comparar a simulação de percurso com medidas externas de potência e
        velocidade em condições controladas;
    \item validar a compilação e o uso em iOS/iPadOS, tablet Android e
        \textit{desktop} macOS;
    \item ampliar a base de rolos, sensores de frequência cardíaca e medidores
        de potência.
\end{itemize}

No eixo de artefato, os próximos passos são:

\begin{itemize}
    \item \textbf{automação do reconhecimento}: detectar subidas, segmentos e
        curvas, marcar pontos críticos e comparar tentativas;
    \item \textbf{análise de pacing}: estimar esforço por trecho e sugerir
        distribuição de potência para a prova;
    \item \textbf{mapas offline}: armazenar tiles e terreno para uso sem rede;
    \item \textbf{calibração e incerteza}: registrar o estado de calibração do
        rolo e propagar incerteza nas métricas;
    \item \textbf{multiplayer local}: permitir sessões conjuntas em rede local,
        sem infraestrutura em nuvem;
    \item \textbf{IA com recuperação}: usar o histórico do atleta como contexto
        para recomendações, mantendo a inferência local.
\end{itemize}

\section{Considerações finais}
\label{sec:consideracoes-finais}

O \argus{} não resolve todos os problemas do ciclismo virtual e não substitui o
reconhecimento presencial, o treinador ou o laboratório. Ele oferece uma
alternativa aberta para quem quer controlar o percurso e o dado, e um artefato
de pesquisa para estudar a integração entre simulação e fisiologia. A
contribuição mais duradoura talvez não seja o código em si, mas a demonstração
de que é possível documentar, de forma reprodutível, o que um software esportivo
mede, como mede e para que serve.
